(1) 珠的位置在惯性系中可取 $r_\perp=a\sin\theta$、$z=-a\cos\theta$，随环旋转时方位角满足 $\dot\phi=\Omega$。因而动能 $T=\tfrac12ma^2(\dot\theta^2+\Omega^2\sin^2\theta)$，重力势以底部为零为 $U_g=mga(1-\cos\theta)$。代入 $L=T-U_g$ 即得只含 $\theta$ 和 $\dot\theta$ 的约化拉格朗日量。
(2) Euler–Lagrange 方程为 $ma^2\ddot\theta=ma^2\Omega^2\sin\theta\cos\theta-mga\sin\theta$，即 $\ddot\theta=\sin\theta(\Omega^2\cos\theta-g/a)$。定义 $U_{\rm eff}=mga(1-\cos\theta)-\tfrac12ma^2\Omega^2\sin^2\theta$，则方程正是 $ma^2\ddot\theta=-dU_{\rm eff}/d\theta$，环外驱动的做功已体现在旋转势里。
(3) 平衡条件 $\sin\theta(g-a\Omega^2\cos\theta)=0$，有底点 $\theta=0$、顶点 $\theta=\pi$，若 $\Omega^2a\ge g$ 另有成对平衡 $\theta=\pm\arccos[g/(a\Omega^2)]$。底点曲率 $U_{\rm eff}''(0)=ma(g-a\Omega^2)$，低于临界稳定；顶点曲率 $-ma(g+a\Omega^2)<0$ 始终不稳定；超临界对称点曲率为 $ma^2\Omega^2\sin^2\theta_0>0$。临界转速为 $\Omega_c=\sqrt{g/a}$。
(4) 由曲率除以 $ma^2$ 得底部振频 $\omega_{\rm low}=\sqrt{g/a-\Omega^2}$，超临界成对平衡频率 $\omega_{\rm high}=\Omega|\sin\theta_0|=\sqrt{\Omega^2-g^2/(a^2\Omega^2)}$。恒速驱动下约化拉格朗日量不显含时间，因此雅可比量 $J=\tfrac12ma^2\dot\theta^2+U_{\rm eff}(\theta)$ 恒定；它一般不等于珠的惯性系机械能，因为旋转外源可交换功。
